% L15a strategy deck: the random-global picking rule.
% One short deck per rule on the strategy menu, read next to the rule's code file
% (strategies/picking/random_global.jl). Written by make_strategy_decks.py (session scratchpad, 2026-11-27) and then
% kept by hand: edit this file, not the generator. Build with make in this folder: LaTeX runs from strategies/,
% so paths are relative to it.
\documentclass[aspectratio=169,11pt]{beamer}
\usepackage{vnslides}

\newcommand{\repo}{https://github.com/varnerlab/CHEME-5660-CourseRepository-Fall-2026/blob/main/lectures/week-15/L15a}

\title{\texorpdfstring{{\vnheavy L15a}}{L15a}: Random Global}
\subtitle{Picking Rule, CHEME 5660 Strategy Menu, Cornell University}
\author{Copyright Jeffrey Varner 2026. All Rights Reserved}

\begin{document}

\vntitlepage

\begin{frame}{Disclaimer and Risks}
  \vspace{0.6em}
  \textbf{\ink{Training only.}} This content is for training and informational
  purposes only. The teaching team does not make, give, or endorse any offer or
  solicitation to buy or sell securities or derivative products, or provide
  investment or trading advice or strategy.

  \vspace{1.1em}
  \textbf{\ink{Risk.}} Trading involves risk and is generally not recommended for
  those with limited resources, experience, or a low risk tolerance. Past
  performance does not guarantee future success.

  \vspace{1.1em}
  \textbf{\ink{Responsibility.}} You are responsible for your investment decisions
  based on your financial circumstances, objectives, risk tolerance, and
  liquidity needs.
\end{frame}

\begin{frame}{The Idea}
  \vspace{0.3em}
  The random-global rule draws 13 tickers at random from the whole universe, the same basket size as the
  global-gamma and correlation-gamma rules:
  \begin{itemize}\setlength{\itemsep}{0.5em}
    \item \hilite{Universe}: 491 stocks and 14 exchange-traded funds (ETFs), without \texttt{SPY} and the Treasury bill fund \texttt{SGOV}.
    \item \hilite{Reason}: a random basket shows how a basket of the same size does by luck alone.
    \item \hilite{Seeds}: the saved results run each random rule with 30 seeds in the backtests and 5 in the stress tests.
  \end{itemize}

  \vspace{0.5em}
  The engine draws a new basket at each decision: on the start day and whenever a trigger fires.
\end{frame}

\begin{frame}{The Rule}
  \vspace{0.3em}
  On decision day $t$:
  \begin{enumerate}\setlength{\itemsep}{0.5em}
    \item Draw $K=13$ distinct tickers at random from the universe, each set of 13 equally likely.
    \item These tickers form the basket $\mathcal{B}(t)$, in the order of the universe.
  \end{enumerate}

  \vspace{0.5em}
  The random stream depends on the run's seed and the decision date, so the same seed and date give the same basket. Unlike the gamma rules, the rule ignores the preference weights.
\end{frame}

\begin{frame}{The Basket of a Made-Up Market}
  \vspace{0.3em}
  In a made-up market of 12 sectors, the rule draws 13 tickers at random:
  \begin{center}
    \includegraphics{slides/figs/Fig-L15a-Deck-Picks-Global.pdf}
  \end{center}
  The global-gamma rule, on the left, takes the 13 tickers with the highest preference weights instead.
\end{frame}

\begin{frame}{Reading the Code}
  \vspace{0.3em}
  The rule's code file is
  \href{\repo/strategies/picking/random_global.jl}{\texttt{strategies/picking/random\_global.jl}}:
  \begin{itemize}\setlength{\itemsep}{0.5em}
    \item \texttt{pick\_random\_global(tickers, rng; count)} draws \texttt{count} tickers from the random stream \texttt{rng}.
    \item \texttt{MyL15aRandomGlobalScreen} is the type the engine calls. It holds the setting \texttt{count~=~13}.
    \item \texttt{MPC.pick(screen, md, t; seed, ...)} builds one random stream for the decision date with the run's \texttt{seed}, runs the rule, and returns the basket's columns in the market data.
  \end{itemize}

  \vspace{0.5em}
  The \href{\repo/strategies/CHEME-5660-L15a-Strategy-Menu-Fall-2026.ipynb}{strategy menu} lists every rule, and the \href{\repo/engine/CHEME-5660-L15a-Trading-Engine-Fall-2026.ipynb}{trading engine} describes the triggers that start each decision.
\end{frame}

\end{document}
